\documentclass[10pt,a4paper]{amsart}
\usepackage[margin=19mm]{geometry}
\usepackage{amsmath,amssymb,mathtools}
\usepackage[scr=rsfs,cal=euler]{mathalfa}
\usepackage{xfrac}
\usepackage[unicode,hidelinks,bookmarks=false]{hyperref}
\newcommand{\ie}{i.e.\ }
\newcommand{\cf}{cf.\ }
\newcommand{\resp}{respectively\ }
\newcommand{\Subsection}{\subsection}
\providecommand{\nobreakdash}{\nobreak\hbox{-}}
\setlength{\emergencystretch}{2em}
\multlinegap0pt
\usepackage{amssymb}
\usepackage{mathtools}
\usepackage[shortlabels]{enumitem}
\usepackage{xcolor}
\usepackage{dsfont}
\newtheorem{assumption}{Assumption}
\newtheorem{lemma}{Lemma}
\usepackage{cleveref}
\crefname{assumption}{Assumption}{Assumptions}
\crefname{lemma}{Lemma}{Lemmas}
\newcommand{\E}{\mathbb E}
\newcommand{\one}{\mathds{1}}
\newcommand{\pro}{\mathbb P}
\title{Multi-Turn LLM Conversations under the Least-Recently-Used Policy: Mean-Field Asymptotics and Hit Ratio Approximation}
\author{Heyuan Yao, Chutong Gao, Yuan Lyu, Izzy Grosof, David Simchi-Levi}
\date{Source: arXiv:2609.02027v1 (CC BY 4.0)}
\begin{document}
\maketitle

\noindent\emph{Source and licence.} Multi-Turn LLM Conversations under the Least-Recently-Used Policy: Mean-Field Asymptotics and Hit Ratio Approximation. Version arXiv:2609.02027v1 (CC BY 4.0), distributed by arXiv under the Creative Commons Attribution 4.0 International licence, http://creativecommons.org/licenses/by/4.0/, as recorded in the arXiv licence metadata for this exact version and shown on the abstract page https://arxiv.org/abs/2609.02027v1. Full licence notice: \texttt{licenses/arxiv-2609.02027-CC-BY-4.0.txt}.\par\smallskip

\noindent\textit{Editorial scope.} Counted unit: the article's lemma Representation (source0.tex lines 913--937). The article's complete original proof of it is delivered with it (source0.tex lines 1157--1246, one \texttt{proof} environment of 90 lines, which the article places away from the statement). No mathematics is written, repaired or re-derived: the statement and the proof are the article's own lines, in the article's own notation, and no retained line is substituted or re-flowed.  Retained uncounted as the source-local context the proof needs: lines 432--435; lines 672--675; lines 677--680; lines 1097--1108.  Nothing was omitted from any retained span, so the package carries no omission record.  No retained line contains a citation, so the wrapper carries no bibliography and no bibliography entry.  No retained line contains a figure environment, an image-inclusion command or drawing code, so no image is delivered and the package declares no asset dependency.  Editorial formatting only, in addition: an AMS \texttt{amsart} preamble that restates the article's own declarations and macros used by the retained lines, namely the environments \texttt{assumption}, \texttt{lemma} on separate counters, together with the standard packages those lines need.  The retained environments print in this document as Assumption 1, Lemma 1 in order of appearance, whereas the article prints the same items with the numbering of its own full document.  The complete native source of the article as distributed in the arXiv e-print for this exact version is bundled unchanged as \texttt{sources/209163/source0.tex}.
\begin{assumption}[i.i.d. Prompt Interarrival Times with a Lipschitz-continuous CDF.]
    \label{ass: Exponential Prompt Inter-arrival Time}
    For each conversation, let $W_j$ denote the interarrival time between the submissions of the turn-$j$ prompt and the turn-$(j+1)$ prompt. The sequence $\{W_j\}_{j\geq1}$ is assumed to be i.i.d. with cumulative distribution function (CDF) $F$, independent of $M$ and the block increments $\{(X_j,Y_j)\}_{j\geq1}$. We further assume that $F$ is Lipschitz-continuous with constant $L_F$ and $\mu_F:= \E [W] <\infty$.
\end{assumption}

    \begin{equation}
    \label{equ: displacement pointwise concentration}
        \pro\left(\left|\frac{D_N(t)}{N}-d(t)\right|>\varepsilon\right)\leq\frac{V_T}{N\varepsilon^2}.
    \end{equation}

    \begin{equation}
    \label{equ: displacement sample path concentration}
        \pro\left(\sup_{0\leq t\leq T}\left|\frac{D_N(t)}{N}-d(t)\right|>\varepsilon\right)\leq\frac{C_T}{N\varepsilon^3},
    \end{equation}

\begin{lemma}[Mean Displacement Curve]
\label{lem: Mean Displacement Curve}
    For every $t\geq0$, the displacement process $D_{N,j}(t)$ has the same law as $D_N(t)$ for every fixed tagged turn state $j$ with $S_j>0$. Moreover,
    \begin{equation}
    \label{equ: Mean Displacement Curve}
        \frac{1}{N}\E[D_N(t)]=d(t):= \lambda_0\E[B]t + \lambda_0A_H\int_0^t\bar F(s)\,ds,
    \end{equation}
    where $\bar F(t):=1-F(t)$ denotes the tail probability of $F$. Moreover, $d$ is strictly increasing and Lipschitz-continuous, with
    \begin{equation*}
        \lambda_0\E[B](t-s) \leq d(t)-d(s) \leq \lambda_0(\E[B]+A_H)(t-s), \qquad\text{for all } 0\leq s\leq t.
    \end{equation*}
\end{lemma}

\begin{lemma}[Poisson Random Measure Representation of the Displacement Process]
\label{lem: Poisson Random Measure Representation}
    Fix a finite horizon $T<\infty$. Recall \cref{ass: Exponential Prompt Inter-arrival Time}, where $\{W_j\}_{j\geq1}$ are i.i.d. interarrival times with cumulative distribution function $F$ and finite mean $\mu_W:=\E[W_1]$, independent of the conversation mark $(M,\{X_j,Y_j\}_{j\geq1})$. Then, for every fixed tagged turn state $j$ with $S_j>0$, the displacement process $\{D_{N,j}(t):0\leq t\leq T\}$ admits the representation
    \begin{equation}
    \label{equ: Poisson Random Measure Representation}
        D_{N,j}(t)\stackrel{\mathcal L}{=}D_N(t):=\int_{\mathcal Z}\Phi_t(z)\,\mathcal N_N(dz), \qquad 0\leq t\leq T,
    \end{equation}
    where $\mathcal N_N$ is a Poisson random measure on the marked conversation trajectory space $\mathcal Z$ with intensity measure $N\nu$. The path functional $\Phi_t(z)$ denotes the cumulative KV block mass contributed by the background conversation $z$ to the displacement process during $(0,t]$.

    For a marked conversation $z=(a,\zeta)$, where $a$ denote its arrival time and $\zeta$ denote its mark, we define its lifetime span $L(\zeta):=\sum_{\ell=1}^{M-1}W_\ell,$ and the envelope
    \begin{equation*}
        G_T(a,\zeta):=B(\zeta)\one\{-L(\zeta)<a\leq T\}.
    \end{equation*}
    Then, for every $z\in\mathcal Z$,
    \begin{equation}
    \label{equ: PRM contribution bound}
        0\leq\Phi_s(z)\leq\Phi_t(z)\leq G_T(z), \qquad 0\leq s\leq t\leq T.
    \end{equation}
    Moreover,
    \begin{equation}
    \label{equ: PRM second moment constant}
        V_T:=\int_{\mathcal Z}G_T(z)^2\,\nu(dz)=\lambda_0T\E[B^2]+\lambda_0\mu_W\E[(M-1)B^2]<\infty.
    \end{equation}
    % Consequently, the law of the background displacement process is invariant with respect to the tagged turn state $j$.
\end{lemma}

\begin{proof}

    For every $t\in[0,T]$, the Poisson random measure representation in \cref{lem: Poisson Random Measure Representation} gives that
    \begin{equation}
    \label{equ: displacement variance exact}
        \text{Var}\left(\frac{D_N(t)}{N}\right)=\frac{1}{N}\int_{\mathcal Z}\Phi_t^2(z)\,\nu(dz) \leq \frac{1}{N}\int_{\mathcal Z}G_T^2(z)\,\nu(dz) =\frac{V_T}{N}.
    \end{equation}
    In addition, due to \cref{lem: Mean Displacement Curve}, $\frac{1}{N}\E[D_N(t)] = d(t)$, we have that 
    \begin{equation}
        \label{equ: displacement MSE upper bound}
        \sup_{0\leq t\leq T}\E\left[\left|\frac{D_N(t)}{N}-d(t)\right|^2\right]\leq\frac{V_T}{N}.
    \end{equation}
    Therefore, applying Chebyshev's inequality, we first derive the following pointwise second-moment bound for any $\varepsilon>0$:
    \begin{equation*}
        \pro\left(\left|\frac{D_N(t)}{N}-d(t)\right|>\varepsilon\right)\leq\frac{1}{\varepsilon^2}\E\left[\left|\frac{D_N(t)}{N}-d(t)\right|^2\right]\leq\frac{V_T}{N\varepsilon^2}.
    \end{equation*}

    We next demonstrate the upper bound for the sample-path $\mathbb L^\infty$-difference. Recall from \cref{lem: Mean Displacement Curve} that for any $t>0$,
    \begin{equation*}
        d'(t)=\lambda_0\E[B]+\lambda_0A_H \bar F(t) \leq L := d'(0)=\lambda_0\E[B]+\lambda_0A_H.
    \end{equation*}
    Hence, the mean displacement curve $d$ is $L$-Lipschitz.

    Fix $0<\varepsilon\leq1$. Choose an equal partition
    \begin{equation*}
        \Pi = \{t_0,\cdots,t_K\},\quad \text{where}\quad 0=t_0<t_1<\cdots<t_K=T,
    \end{equation*}
    satisfying that
    \begin{equation*}
    \|\Pi\| : =\max_{0\leq k<K}(t_{k+1}-t_k)\leq\frac{\varepsilon}{4 L }.
    \end{equation*}
    It is enough to choose $K+1\leq 2+\frac{4 L T}{\varepsilon}$, i.e., $K=  \lfloor1+\frac{4 L T}{\varepsilon}\rfloor$.

    Fix $t\in[t_k,t_{k+1}]$. Because $D_N(t)$ and  $d(t)$ are both nondecreasing,
    \begin{equation*}
        D_N(t_k)\leq D_N(t)\leq D_N(t_{k+1}) \quad \text{and} \quad d(t_k)\leq d(t)\leq d(t_{k+1}).
    \end{equation*}
    Therefore, 
    \begin{align*}
        \frac{D_N(t)}{N}-d(t)
        &\leq\frac{D_N(t_{k+1})}{N}-d(t_k) \\
        &=\left(\frac{D_N(t_{k+1})}{N}-d(t_{k+1})\right)+\left(d(t_{k+1})-d(t_k)\right) \\
        &\leq\left|\frac{D_N(t_{k+1})}{N}-d(t_{k+1})\right|+\frac{\varepsilon}{4}.
    \end{align*}
    and similarly,
    \[
    d(t)-\frac{D_N(t)}{N} \leq \left|\frac{D_N(t_{k})}{N}-d(t_{k})\right|+\frac{\varepsilon}{4}.
    \]

    To establish 
    \begin{equation*}
        \sup_{0\leq t\leq T}\left|\frac{D_N(t)}{N}-d(t)\right|\leq\varepsilon,
    \end{equation*}
    it is enough to let, for any $1\leq k\leq K$ 
    \begin{equation*}
        \max_{0\leq k\leq K}\left|\frac{D_N(t_k)}{N}-d(t_k)\right|\leq\frac{3\varepsilon}{4}.
    \end{equation*}
    Therefore, we can bound the $L^\infty$-distance via a union bound, such that
    \begin{align*}
        \pro\left(\sup_{0\leq t\leq T}\left|\frac{D_N(t)}{N}-d(t)\right|>\varepsilon\right)& \leq\sum_{k=0}^{K}\pro\left(\left|\frac{D_N(t_k)}{N}-d(t_k)\right|>\frac{3\varepsilon}{4}\right)\\
        & \leq (K+1) \frac{16V_T}{9N\varepsilon^2}\\
        & \leq (2+\frac{4 L T}{\varepsilon})\frac{16V_T}{9N\varepsilon^2}.
    \end{align*}

    When $0<\varepsilon\leq1$, the last term can be further bounded by  
    \begin{equation*}
        (2+\frac{4 L T}{\varepsilon})\frac{16V_T}{9N\varepsilon^2} \leq \frac{2+4 L T}{\varepsilon}\frac{16V_T}{9N\varepsilon^2} = \frac{C_T}{N\varepsilon^3}.
    \end{equation*}  


    We finally show the convergence result with the given rate. For the pointwise convergence rate, set $\varepsilon = \delta/\sqrt N$ for some fixed $\delta>0$ in \eqref{equ: displacement pointwise concentration}. Then
    \begin{equation*}
        \pro\left(\sqrt N\left|\frac{D_N(t)}{N}-d(t)\right|>\delta \right)\leq\frac{V_T}{\delta^2}.
    \end{equation*}
    Taking $\limsup_{N\to\infty}$ and then $\delta\to\infty$, we have that
    \begin{equation*}
        \left|\frac{D_N(t)}{N}-d(t)\right|=O_{\mathbb P}(N^{-1/2}).
    \end{equation*}
    
    For the sample-path convergence rate, set $\varepsilon=\delta N^{-1/3}$ in \eqref{equ: displacement sample path concentration}. For every fixed $\delta>0$ and all sufficiently large $N$, $\varepsilon\leq 1$, and
    \begin{equation*}
        \pro\left(N^{1/3}\sup_{0\leq t\leq T}\left|\frac{D_N(t)}{N}-d(t)\right|>\delta\right)\leq\frac{C_T}{\delta^3}.
    \end{equation*}
    Taking $\limsup_{N\to\infty}$ and then $\delta\to\infty$, we have that
    \begin{equation*}
        \sup_{0\leq t\leq T}\left|\frac{D_N(t)}{N}-d(t)\right|=O_{\mathbb P}(N^{-1/3}).
    \end{equation*}

    Finally, because $D_{N,j}(\cdot)\stackrel{\mathcal L}{=}D_N(\cdot)$ for every fixed turn state $j$ with $S_j>0$, the same concentration bounds, convergence in probability, and finite-size rates hold for the normalized tagged-conversation displacement process $\frac{D_{N,j}(\cdot)}{N}$. This completes the proof.
\end{proof}

\end{document}
